\begin{table}[htbp]
  \centering
  \caption{Mean time for one simulation step: the scene total, median over repeats, divided by the number of frames. A step runs both query types, so this is the cost of the vertex-face and edge-edge work of that frame together. \emph{BP full} is the whole broad phase, the acceleration structure plus the queries over it; \emph{BP prep} is the structure alone and \emph{BP queries} the queries alone. The narrow phase is named by the answer it is asked for: \emph{NP EToI} returns one earliest time of impact for the step, so every query prunes against the running minimum, and \emph{NP per-pair} returns one per candidate with no shared bound. BP full here includes building the swept boxes.}
  \label{tab:per-frame-relaxed}
  \fittable{%
  \begin{tabular}{lrlrrr}
    \toprule
    scene & frames & mode & BP full (ms) & NP EToI (ms) & total ms \\
    \midrule
    armadillo-rollers & 396 & GPU & 3.01 & 3.29 & 6.30 \\
    armadillo-rollers & 396 & CPU & 3.53 & 1.66 & 5.19 \\
    cloth-ball & 43 & GPU & 15.02 & 3.94 & 18.96 \\
    cloth-ball & 43 & CPU & 13.66 & 23.37 & 37.03 \\
    cloth-funnel & 372 & GPU & 4.17 & 1.70 & 5.88 \\
    cloth-funnel & 372 & CPU & 2.36 & 0.54 & 2.90 \\
    n-body & 74 & GPU & 53.63 & 6.34 & 59.98 \\
    n-body & 74 & CPU & 43.70 & 121.30 & 165.01 \\
    puffer-ball & 120 & GPU & 178.75 & 6.22 & 184.97 \\
    puffer-ball & 120 & CPU & 209.63 & 269.55 & 479.19 \\
    rod-twist & 2,556 & GPU & 3.40 & 3.49 & 6.89 \\
    rod-twist & 2,556 & CPU & 6.39 & 13.89 & 20.28 \\
    \bottomrule
  \end{tabular}}
  \par\smallskip\footnotesize A mean rather than a median over steps: the scene total is what a run costs, and the mean is the only average that divides back into it.
  \par\smallskip\footnotesize Source: \texttt{benchmark/assessment/broadphase-cell2dmin.csv}
\end{table}
